\section{Constructing a negative kernel correlation}
\label{sec:kernel}

This section reduces Theorem~\ref{thm:main} to a finite-field counting
lemma. We first state that lemma, then construct a nonnegative kernel
whose homomorphism density is too small. All probability spaces in
the construction are finite; their atom probabilities need not be
rational.

\subsection{The rank-layer model}

Let $q$ be an odd prime, let $\chi:\F_q\to\{0,1,-1\}$ be the
quadratic character with $\chi(0)=0$, and let $\mathcal S_D$ be the
space of symmetric $D$-by-$D$ matrices over $\F_q$, equipped with
uniform probability. For a symmetric matrix of positive rank, its
\emph{sign} is the quadratic character of the determinant of its
induced nondegenerate symmetric form on the quotient by the radical.
A change of basis multiplies this determinant by a square, so the
sign is well defined. Fix a positive even dimension $D=2r$ and set
\begin{equation}
 K_\xi(Z)=
 \frac{\1_{\{\rank Z=r,\ \operatorname{sign}(Z)=\xi\}}}
 {\P_{Z'\in\mathcal S_D}(\rank Z'=r,\ \operatorname{sign}(Z')=\xi)},
 \qquad \xi\in\{1,-1\},\qquad c_0=\chi((-1)^r).
 \label{eq:kernel-rank-definition}
\end{equation}
Both denominators are positive: diagonal matrices realize either
nonzero determinant square class in rank $r$. In particular,
$K_\xi\ge0$ and $\E K_\xi=1$.

\begin{lemma}[Rank-layer limit]
\label{lem:rank-layer}
There is an integer $D_0$ such that, for every fixed even $D\ge D_0$,
the following assertion holds as $q$ tends to infinity through odd
primes. Let $X_i$, $i\in I$, and $Y_j$, $j\in J$, be independent
uniform elements of $\mathcal S_D$. For every set
$B\subseteq\mathscr C$ and every choice of signs
$\xi_{ij}\in\{1,-1\}$ on $B$,
\begin{equation}
 \E\prod_{(i,j)\in B}K_{\xi_{ij}}(X_i-Y_j)
 =\E_\eta\prod_{j\in J}
 \prod_{\substack{e=\{i,k\}\subset j\\(i,j),(k,j)\in B}}
 \bigl(1+c_0\eta_e\xi_{ij}\xi_{kj}\bigr)+o(1),
 \label{eq:rank-layer}
\end{equation}
where the $\eta_e$, $e\in E$, are independent uniform signs. The error
is uniform over all such $B$ and sign choices. The same assertion
holds with $Y_j-X_i$ in place of $X_i-Y_j$.
\end{lemma}

The averaging on the right has a useful exact form. For
\(e=\{v,w\}\in E\), let \(j,l\) be its two faces and put
\begin{equation}
 U_e=\{(v,j),(w,j),(v,l),(w,l)\}\subseteq\mathscr C.
 \label{eq:kernel-four-corners}
\end{equation}
Write \(\xi_\kappa=\xi_{ij}\) for \(\kappa=(i,j)\in B\).
If only one of the two face factors on \(e\) is present, its average
over \(\eta_e\) is \(1\). If both are present, their average is
\(1+\prod_{\kappa\in U_e}\xi_\kappa\), since \(c_0^2=1\).
Independence of the \(\eta_e\) therefore gives
\begin{equation}
 \E_\eta\prod_{j\in J}
 \prod_{\substack{e=\{i,k\}\subset j\\(i,j),(k,j)\in B}}
 (1+c_0\eta_e\xi_{ij}\xi_{kj})
 =
 \prod_{\substack{e\in E\\U_e\subseteq B}}
 \left(1+\prod_{\kappa\in U_e}\xi_\kappa\right).
 \label{eq:kernel-parity-model}
\end{equation}
Thus every fully active pair tests the parity of its four incidence
signs, contributing \(0\) or \(2\). This is an identity for the
integrated sign model in Equation~\eqref{eq:rank-layer}. The
quantification over \(B\) is needed because each coordinate in the
activation construction below selects its own active incidences.

The proof of Lemma~\ref{lem:rank-layer} occupies
Sections~\ref{sec:transverse}--\ref{sec:tail}.
The dimension is fixed before the prime tends to infinity; uniformity
in $B$ and the signs then concerns a fixed finite family. For the
assertion with reversed matrix differences, negation multiplies the sign of a rank-$r$ matrix by
$c_0$ and bijects the corresponding layers, whence
$K_\xi(-Z)=K_{c_0\xi}(Z)$. Replacing every incidence sign by
$c_0\xi_{ij}$ leaves each sign product on the right of
Equation~\eqref{eq:rank-layer} unchanged.

For a kernel $W:\mathcal X\times\mathcal Y\to[0,\infty)$ on two
finite probability spaces, write
\begin{equation}
 t_{I,J}(W)=\E\prod_{(i,j)\in\mathscr C}W(x_i,y_j),
 \qquad W^\top(y,x)=W(x,y),
 \label{eq:kernel-bipartite-density}
\end{equation}
where all $x_i\in\mathcal X$ and $y_j\in\mathcal Y$ are sampled
independently. Thus $t_{I,J}(W^\top)$ samples the $I$-vertices from
$\mathcal Y$ and the $J$-vertices from $\mathcal X$.

For a symmetric kernel $U:\Omega\times\Omega\to[0,\infty)$ on one
finite probability space, write $t(H,U)=t_{I,J}(U)$, with every vertex
sampled independently from $\Omega$; sampled labels may coincide.

\subsection{Activation laws and the coordinate expansion}

A type records which point or face of the fixed complex a sampled
vertex is intended to represent. Independent matrix coordinates carry
separate copies of the rank-layer model, and an activation law chooses
which copies occur at a pair of types. We first leave these finite
choices arbitrary and derive their contribution to the density.

Let $\mathcal A,\mathcal B$ be finite type sets, with probability
distributions $\pi,\nu$, and let $\mathcal K$ be a fixed finite
coordinate set. Use the spaces
\[
 \mathcal X=\mathcal A\times\mathcal S_D^{\mathcal K},\qquad
 \mathcal Y=\mathcal B\times\mathcal S_D^{\mathcal K},
\]
with the type distributions just specified and all matrix coordinates
independent and uniform. For each $(a,b)\in\mathcal A\times\mathcal B$,
choose a probability law on a subset $S\subseteq\mathcal K$ and signs
$(\xi_\kappa)_{\kappa\in S}$. Define
\begin{equation}
 W_q\bigl((a,x),(b,y)\bigr)
 =\E_{S,\xi\mid a,b}\prod_{\kappa\in S}
 K_{\xi_\kappa}(x^\kappa-y^\kappa).
 \label{eq:kernel-activation}
\end{equation}
An empty product is one. For fixed types, averaging over the matrices
gives one, since each difference is uniform and distinct coordinates
are independent. Therefore $\E W_q=1$, for every $q$.

For $\alpha\subseteq\mathcal K$ put
\[
 \phi_\alpha(S,\xi)=
 \1_{\{\alpha\subseteq S\}}\prod_{\kappa\in\alpha}\xi_\kappa,
 \qquad
 A_{ab}(\alpha,\beta)=\E_{S,\xi\mid a,b}\phi_\alpha\phi_\beta.
\]
Here $\phi_\alpha$ is defined to be zero when $\alpha\not\subseteq S$;
in particular $\phi_\varnothing=1$. We will impose the moment condition
\begin{equation}
 \E_{S,\xi\mid a,b}\phi_\alpha=0
 \quad\text{for every $a,b$ and every nonempty $\alpha$.}
 \label{eq:kernel-zero-moments}
\end{equation}
If $\alpha\cap\beta=\varnothing$, then
$\phi_\alpha\phi_\beta=\phi_{\alpha\cup\beta}$, so
\begin{equation}
 A_{ab}(\alpha,\beta)=0
 \quad\text{if $\alpha\cap\beta=\varnothing$ and
 $\alpha\cup\beta\ne\varnothing$.}
 \label{eq:kernel-disjoint-zero}
\end{equation}

At a corner $(i,j)$, let $e(i,j),f(i,j)$ denote the two point pairs
of $j$ containing $i$; their order is immaterial. Expanding
Equation~\eqref{eq:kernel-parity-model} at every coordinate selects,
for each pair \(e\), a subset \(\alpha_e\subseteq\mathcal K\) of
coordinates at which its four-sign term is used. At a corner, the
two labels meeting there contribute their overlap moment
\(A_{ab}(\alpha_{e(i,j)},\alpha_{f(i,j)})\).
For type assignments
$a:I\to\mathcal A$, $b:J\to\mathcal B$, define
\begin{equation}
 C_A(a,b)=
 \sum_{\substack{\alpha_e\subseteq\mathcal K\ (e\in E)\\
                   \alpha_e\ne\varnothing\ (e\in E)}}
 \prod_{(i,j)\in\mathscr C}
 A_{a(i)b(j)}\bigl(\alpha_{e(i,j)},\alpha_{f(i,j)}\bigr).
 \label{eq:kernel-coefficient}
\end{equation}
These are finite real numbers, with no restriction on either type
map. In particular, repetitions of types are included.
The next lemma identifies this all-nonempty label sum as the
limiting nonconstant contribution at the fixed type assignment.

\begin{lemma}[Coordinate expansion]
\label{lem:kernel-expansion}
Assume Equation~\eqref{eq:kernel-zero-moments} and fix an even dimension
as in Lemma~\ref{lem:rank-layer}. For fixed type laws and activation
laws,
\begin{align}
 t_{I,J}(W_q)&=1+
 \sum_{a:I\to\mathcal A\,,\ b:J\to\mathcal B}
 \left(\prod_{i\in I}\pi_{a(i)}\right)
 \left(\prod_{j\in J}\nu_{b(j)}\right)C_A(a,b)+o(1).
 \label{eq:kernel-expansion}
\end{align}
The corresponding formula for $W_q^\top$ has type maps
$b:I\to\mathcal B$, $a:J\to\mathcal A$ and the corner factors
$A_{a(j)b(i)}$.
\end{lemma}

\begin{proof}
Conditional on all types, sample an activation law independently at
each incidence when expanding the product in
Equation~\eqref{eq:kernel-bipartite-density}. Conditional on these
samples, the matrix integral factors across coordinates. At coordinate
$\kappa$ apply Lemma~\ref{lem:rank-layer} with the active set of
incidences for which $\kappa\in S$, and with their sampled signs.
Use Equation~\eqref{eq:kernel-parity-model} at each coordinate.
There are only finitely many activation and sign configurations.
The model expectations are bounded, and their approximating matrix
expectations are consequently bounded for large $q$. A finite product
and then the finite average therefore preserve the $o(1)$ error.

Expand these finite parity products, and let \(\alpha_e\) be the set
of coordinates in which the four-sign term for \(e\) is selected. A
coordinate selected for \(e\) requires activation at all four
corners in \(U_e\) and contributes the product of their signs.
At $(i,j)$ the two subsets on its incident pairs give exactly
$\phi_{\alpha_{e(i,j)}}\phi_{\alpha_{f(i,j)}}$; averaging that
incidence's activation and signs gives its factor \(A\). Thus the
limiting contribution at the fixed type assignment is
\begin{equation}
 \sum_{\alpha_e\subseteq\mathcal K\ (e\in E)}
 \prod_{(i,j)\in\mathscr C}
 A_{a(i)b(j)}
   \bigl(\alpha_{e(i,j)},\alpha_{f(i,j)}\bigr).
 \label{eq:kernel-unrestricted-labels}
\end{equation}
Repeated vertex types leave the matrix draws at distinct source vertices
independent. Repeated type pairs use independent samples from the same
activation law at their distinct source incidences.

The all-empty labeling contributes one. In any other nonzero term,
no face can have both an empty and a nonempty pair label: the two
meet at a corner, where Equation~\eqref{eq:kernel-disjoint-zero} makes
the factor zero. Hence any face with one nonempty label has all three
labels nonempty. The shared-pair labels and connectedness of the
face-neighbor graph propagate this conclusion to every face. Thus
the remaining labels are exactly those summed in
Equation~\eqref{eq:kernel-coefficient}. This proves
Equation~\eqref{eq:kernel-expansion}. For the transpose, the same
argument uses the last assertion of Lemma~\ref{lem:rank-layer}.
\end{proof}

\subsection{A single negative term}

The problem is now finite: choose activation laws so that the weighted
coefficient sum in Equation~\eqref{eq:kernel-expansion} is negative.
We first make the coefficient at the identity type assignment negative,
and then choose the priors and thinning factors to make it dominate
the other assignments.

Take $\mathcal A=I$, $\mathcal B=J$, and one coordinate for each
corner: $\mathcal K=\mathscr C$. The four-incidence set \(U_e\) from
Equation~\eqref{eq:kernel-four-corners} is now also a subset of the
coordinate set.
At a corner $\kappa=(v,j)$, write $e,f$ for the two pairs of $j$
through $v$ and put
\begin{equation}
 M_\kappa=U_e\cup U_f,
 \qquad \Gamma_\kappa=U_e\mathbin\triangle U_f.
 \label{eq:kernel-seven-six}
\end{equation}
Distinct faces cannot share two distinct pairs. Hence the opposite
faces on $e,f$ are distinct and
$U_e\cap U_f=\{\kappa\}$. In particular $M_\kappa$ has seven
coordinates and $\Gamma_\kappa=M_\kappa\setminus\{\kappa\}$ has six.
Figure~\ref{fig:kernel-corner-activation} displays these sets.

\begin{figure}[htbp]
\centering
\begin{tikzpicture}[font=\small]
  \fill[blue!7] (-1.25,0) ellipse (2.05 and 1.45);
  \fill[orange!9] (1.25,0) ellipse (2.05 and 1.45);
  \draw[blue!60!black,thick] (-1.25,0) ellipse (2.05 and 1.45);
  \draw[orange!70!black,thick] (1.25,0) ellipse (2.05 and 1.45);
  \node[blue!60!black] at (-2.4,1.35) {$U_e$};
  \node[orange!70!black] at (2.4,1.35) {$U_f$};
  \node at (-1.7,.75) {$(w,j)$};
  \node at (-2.0,-.1) {$(v,l)$};
  \node at (-1.45,-.85) {$(w,l)$};
  \node at (1.7,.75) {$(z,j)$};
  \node at (2.0,-.1) {$(v,m)$};
  \node at (1.45,-.85) {$(z,m)$};
  \fill[linkblue] (0,.55) circle (1.7pt);
  \node[align=center] at (0,-.05) {$\kappa$\\$=(v,j)$};
  \node at (0,-1.9) {$M_\kappa=U_e\cup U_f,\qquad
                    \Gamma_\kappa=M_\kappa\setminus\{\kappa\}$};
\end{tikzpicture}
\caption{The seven coordinates available at the corner
$\kappa=(v,j)$, where $j=\{v,w,z\}$, $e=\{v,w\}$ and
$f=\{v,z\}$. The opposite faces on $e$ and $f$ are $l$ and $m$.
The two activation branches keep the same six outer coordinates;
one also keeps the common corner marked by the dot.}
\label{fig:kernel-corner-activation}
\end{figure}


Fix $0<\theta<1$. Give each corner a sign $\sigma_\kappa$, choosing
exactly one negative sign, for example $\sigma_{(0,0)}=-1$ and all
others $+1$. At the type pair $\kappa=(v,j)$ with $v\in j$ use this
\emph{base law}:
\begin{itemize}
\item with probability $1/2$, set $S=M_\kappa$ and give each sign
vector on $S$ probability
$2^{-7}(1+\sigma_\kappa\theta\prod_{\gamma\in\Gamma_\kappa}\xi_\gamma)$;
\item with probability $1/2$, set $S=\Gamma_\kappa$ and give each
sign vector on $S$ probability
$2^{-6}(1-\sigma_\kappa\theta\prod_{\gamma\in\Gamma_\kappa}\xi_\gamma)$.
\end{itemize}
These conditional masses are positive and sum to one because a
nonempty sign monomial has zero average under uniform signs. At a
nonincident type pair $(v,j)$ use $S=\varnothing$ deterministically.
Write $A^{\mathrm{base}}$ for the resulting moments.

Under either branch at a corner, the only nonempty sign monomial with
possibly nonzero expectation has index $\Gamma_\kappa$. Its conditional
expectations are respectively $\sigma_\kappa\theta$ and
$-\sigma_\kappa\theta$. The two effects of the branch choice are
\begin{equation}
 \E\phi_{\Gamma_\kappa}
 =\tfrac12\sigma_\kappa\theta-\tfrac12\sigma_\kappa\theta=0,
 \qquad
 A^{\mathrm{base}}_\kappa(U_e,U_f)
 =\tfrac12\sigma_\kappa\theta.
 \label{eq:kernel-branch-biases}
\end{equation}
The second identity holds because only the \(M_\kappa\) branch
contains \(U_e\cup U_f\), while their sign product is the
\(\Gamma_\kappa\) monomial. All base laws satisfy
Equation~\eqref{eq:kernel-zero-moments}, but this overlap moment can
have either sign.
More explicitly, for $Q=\alpha\cup\beta$ and
$R=\alpha\mathbin\triangle\beta$,
\begin{align}
 A^{\mathrm{base}}_\kappa(\alpha,\beta)
 &=\tfrac12\1_{\{Q\subseteq M_\kappa\}}
       \bigl(\1_{\{R=\varnothing\}}
                +\sigma_\kappa\theta\1_{\{R=\Gamma_\kappa\}}\bigr)
 \notag\\
 &\quad+\tfrac12\1_{\{Q\subseteq\Gamma_\kappa\}}
       \bigl(\1_{\{R=\varnothing\}}
                -\sigma_\kappa\theta\1_{\{R=\Gamma_\kappa\}}\bigr).
 \label{eq:kernel-base-moment}
\end{align}

\begin{lemma}[Support forcing]
\label{lem:kernel-support}
At the identity type assignments $a(i)=i$, $b(j)=j$, the sum in
Equation~\eqref{eq:kernel-coefficient} for the base laws has exactly one
nonzero term. Its labels and value are
\begin{equation}
 \alpha_e=U_e\quad(e\in E),\qquad
 C_{A^{\mathrm{base}}}(\mathrm{id},\mathrm{id})
 =-(\theta/2)^{66}.
 \label{eq:kernel-negative-term}
\end{equation}
\end{lemma}

\begin{proof}
We first prove the support identity
\begin{equation}
 \bigcap_{\kappa\in U_g}M_\kappa=U_g
 \qquad(g\in E).
 \label{eq:kernel-support-intersection}
\end{equation}
The same set \(\mathscr C\) indexes coordinates \(\gamma\) and
activation corners \(\kappa\). Availability is symmetric:
\(\gamma\in M_\kappa\) if and only if \(\kappa\in M_\gamma\).
Both conditions mean that the two corners lie in a common set \(U_e\).
Fix a coordinate $\gamma=(v,j)$, write $j=\{v,w,z\}$, and let
$e=\{v,w\}$, $f=\{v,z\}$ have opposite faces $l,m$, respectively.
This coordinate belongs to precisely $U_e$ and $U_f$. Consequently it
is available at exactly the following activation corners: all three
corners of $j$, the two $e$-endpoint corners of $l$, and the two
$f$-endpoint corners of $m$. To be available at all four corners of
$g$, its two faces must therefore be among $j,l,m$. The choice $j,l$
forces $g=e$, and $j,m$ forces $g=f$, since distinct triples share at
most one pair. The choice $l,m$ cannot work: its two endpoints would
both have to lie in $\{v,w\}\cap\{v,z\}=\{v\}$. Thus the coordinate
$\gamma$ is available at every corner of $g$ if and only if
$g\in\{e,f\}$, equivalently if and only if $\gamma\in U_g$.
This proves Equation~\eqref{eq:kernel-support-intersection}.

In a nonzero identity-type term, $\alpha_g$ must lie in the maximal
activation support at each of the four corners of $g$. Hence
$\alpha_g\subseteq U_g$. At a corner $\kappa$ with pairs $e,f$, the
two nonempty labels cannot be disjoint by
Equation~\eqref{eq:kernel-disjoint-zero}. Since
$U_e\cap U_f=\{\kappa\}$, both labels contain $\kappa$. Applying
this at all four corners of $g$ gives $\alpha_g=U_g$.

For these labels, Equation~\eqref{eq:kernel-branch-biases} gives
the local factor \(\sigma_\kappa\theta/2\). Multiplying the $66$ factors,
with exactly one negative $\sigma_\kappa$, proves the assertion.
\end{proof}

Each local matrix $A^{\mathrm{base}}_\kappa$ is positive semidefinite,
being a matrix of second moments of the random variables $\phi_\alpha$.
For example its principal block on $U_e,U_f$ is
\[
 \begin{pmatrix}1/2&\sigma_\kappa\theta/2\\
                 \sigma_\kappa\theta/2&1/2\end{pmatrix},
\]
with eigenvalues $(1+\theta)/2$ and $(1-\theta)/2$. Thus the negative
global term uses an admissible negative off-diagonal moment. The
support argument proves that local diagonal alternatives cannot form
any additional nonempty global labeling at the identity types.

\subsection{Pinning the types by color refinement}

The identity calculation alone does not control the kernel density:
the expansion averages over every type assignment, including maps with
repetitions. We now give the identity strictly larger score than every
competitor. The following lemma applies to any bipartite graph whose
part colors refine to distinct vertex colors.

\begin{lemma}[Pinning by vertex and pair scores]
\label{lem:pinning}
Let $F=(P\sqcup Q,L)$ be a finite bipartite graph for which color
refinement, starting with the two part colors, becomes discrete within
each part. There are real vertex scores $s_t$, $t\in P\sqcup Q$,
and pair scores $s_{ab}$, $(a,b)\in P\times Q$, such that the identity
is the unique maximizer of
\begin{equation}
 \mathcal Q(f)=\sum_{u\in P\sqcup Q}s_{f(u)}
              +\sum_{\{u,v\}\in L,\ u\in P}s_{f(u)f(v)}
 \label{eq:kernel-pinning-score}
\end{equation}
over \emph{all} maps $f$ preserving the two parts. In particular the
maximum has a positive gap from every other such map; no injectivity
or adjacency-preservation condition is imposed.
\end{lemma}

\begin{proof}
Let $\ell_0$ be the two part colors and, at round $t$, let $h_t(u)$
be the histogram of $\ell_{t-1}$-colors among the neighbors of $u$.
The new color is $\ell_t(u)=(\ell_{t-1}(u),h_t(u))$. Choose $R\ge1$
rounds giving a discrete coloring. At round $t$ assign the vertex score
$-\|h_t(a)\|^2$ to type $a$ and the pair score
\[
 2h_t(a)(\ell_{t-1}(b))+2h_t(b)(\ell_{t-1}(a))
\]
to every $(a,b)\in P\times Q$, including nonedges. The total round
score of \(f\) is
\[
 F_t(f)=-\sum_u\|h_t(f(u))\|^2
 +2\sum_u\sum_{v\in N_F(u)}
     h_t(f(u))\bigl(\ell_{t-1}(f(v))\bigr).
\]
Here \(N_F(u)\) is the neighbor set of \(u\). Both sums are indexed
by the vertices and directed incidences of the source graph \(F\),
even when \(f\) repeats a type.

If $\ell_{t-1}(f(u))=\ell_{t-1}(u)$ for every $u$, the inner
neighbor sum equals $\langle h_t(f(u)),h_t(u)\rangle$. Thus
\begin{align}
 F_t(f)
 &=\sum_u\bigl(2\langle h_t(f(u)),h_t(u)\rangle
                   -\|h_t(f(u))\|^2\bigr)\notag\\
 &=\sum_u\|h_t(u)\|^2
       -\sum_u\|h_t(f(u))-h_t(u)\|^2.
 \label{eq:kernel-pinning-square}
\end{align}
The first equality uses only preservation of previous colors. It
therefore applies to noninjective maps and to maps taking edges to
nonedges. In this situation $F_t(\mathrm{id})-F_t(f)$ is a nonnegative
integer, zero precisely when $f$ also preserves all new colors.

To combine the rounds into a single score, set
\[
 M=\max_{f,\,1\le t\le R}|F_t(f)-F_t(\mathrm{id})|,
 \qquad L_0=M+2,
 \qquad \mathcal Q(f)=\sum_{t=1}^R L_0^{R-t}F_t(f).
\]
The maximum defining $M$ is finite. For a nonidentity map, let $t$
be the first round at which a color fails to be preserved. All
earlier scores tie the identity; the loss at round $t$ is at least
$L_0^{R-t}$. Its possible later advantage is at most
\[
 M\sum_{j=0}^{R-t-1}L_0^j<L_0^{R-t}.
\]
Hence its total score is strictly smaller. A nonidentity map must
have such a first round because the final colors are discrete.
The weighted sums of the round vertex scores and pair scores are the
required $s_t,s_{ab}$. Finiteness of the set of maps gives the positive
gap.
\end{proof}

Apply Lemma~\ref{lem:pinning} to $H$, using
Proposition~\ref{prop:complex}. Let its scores be $s_a$ on $I$, $s_b$
on $J$, and $s_{ab}$ on $I\times J$. For a positive parameter $\rho$,
choose priors and thinning probabilities
\begin{equation}
 \pi_a=\frac{e^{\rho s_a}}{Z_I},\qquad
 \nu_b=\frac{e^{\rho s_b}}{Z_J},\qquad
 \lambda_{ab}=e^{\rho(s_{ab}-s_{\max})},\qquad
 s_{\max}=\max_{a,b}s_{ab},
 \label{eq:kernel-priors}
\end{equation}
where $Z_I=\sum_{a\in I}e^{\rho s_a}$ and
$Z_J=\sum_{b\in J}e^{\rho s_b}$. Since $0<\lambda_{ab}\le1$, the
following is a probability law: use the entire base law with
probability $\lambda_{ab}$, and use $S=\varnothing$ otherwise.
Its moments satisfy the exact identity
\begin{equation}
 A^\lambda_{ab}(\alpha,\beta)
 =\lambda_{ab}A^{\mathrm{base}}_{ab}(\alpha,\beta)
 +(1-\lambda_{ab})\1_{\{\alpha=\beta=\varnothing\}}.
 \label{eq:kernel-thinning}
\end{equation}
This preserves Equation~\eqref{eq:kernel-zero-moments}. In every
all-nonempty term it introduces one factor $\lambda_{ab}$ per
incidence, independently of the sizes of the coordinate subsets.

Let $C(a,b)=C_{A^{\mathrm{base}}}(a,b)$ and let $t_0$ be the
coefficient of the nonempty terms in
Equation~\eqref{eq:kernel-expansion} for the thinned laws and these
priors. By Equation~\eqref{eq:kernel-thinning},
\begin{align}
 t_0
 &=Z_I^{-13}Z_J^{-22}e^{-66\rho s_{\max}}
     \sum_{a:I\to I\,,\ b:J\to J}
           C(a,b)e^{\rho\mathcal Q(a,b)}.
 \label{eq:kernel-pinned-coefficient}
\end{align}
All normalizing factors outside the sum are common to every map.
Here $\mathcal Q(a,b)$ means the score of the part-preserving map
specified by $a,b$.

Write $Q_*=\mathcal Q(\mathrm{id},\mathrm{id})$, and let
\[
 \delta_*=\min_{(a,b)\ne(\mathrm{id},\mathrm{id})}
          (Q_*-\mathcal Q(a,b))>0,
 \qquad
 B_*=\sum_{(a,b)\ne(\mathrm{id},\mathrm{id})}|C(a,b)|<\infty.
\]
After division by $e^{\rho Q_*}$, the absolute value of the
competitor sum is at most $B_*e^{-\rho\delta_*}$. Choose a fixed
finite $\rho>0$ large enough that
\begin{equation}
 B_*e^{-\rho\delta_*}<(\theta/2)^{66};
 \label{eq:kernel-finite-gap}
\end{equation}
if $B_*=0$, any positive $\rho$ suffices. By
Equation~\eqref{eq:kernel-negative-term} this makes $t_0<0$.

\subsection{Dilution, symmetrization, and finite hosts}

We have made the limiting nonempty coefficient negative in one
orientation. The symmetric kernel below multiplies the densities in
both orientations. Dilution makes the negative term dominate the
unknown term from the other orientation.

Keep $\theta,\rho$, the priors, and all the thinned laws fixed. Add
an inactive type $*$ to $\mathcal A=I$, giving it probability
$1-\varepsilon$ and giving each old type $a$ probability
$\varepsilon\pi_a$, where $0<\varepsilon<1$. Every pair involving
$*$ uses $S=\varnothing$. Keep the $J$-type prior $\nu$ unchanged.
Let $W_{q,\varepsilon}$ be the resulting kernel.

For precision, define also the finite coefficient for the opposite
orientation before dilution:
\begin{align}
 t_1={}&\sum_{b:I\to J\,,\ a:J\to I}
       \left(\prod_{i\in I}\nu_{b(i)}\right)
       \left(\prod_{j\in J}\pi_{a(j)}\right)
       \sum_{\substack{\alpha_e\subseteq\mathscr C\ (e\in E)\\
                       \alpha_e\ne\varnothing\ (e\in E)}}
       \prod_{(i,j)\in\mathscr C}
       A^\lambda_{a(j)b(i)}
         \bigl(\alpha_{e(i,j)},\alpha_{f(i,j)}\bigr).
 \label{eq:kernel-transpose-coefficient}
\end{align}
Both $t_0<0$ and $t_1\in\R$ are independent of $q$ and
$\varepsilon$.

In a nonempty-label term, a vertex with inactive type gives a zero
factor on each of its incidences. Every vertex of $H$ has positive
degree, so every original $I$-type draw must be active. There are
$13$ such draws for $W_{q,\varepsilon}$ and $22$ for its transpose.
Lemma~\ref{lem:kernel-expansion} therefore gives the exact limiting
coefficients
\begin{align}
 t_{I,J}(W_{q,\varepsilon})
   &=1+\varepsilon^{13}t_0+o(1),\notag\\
 t_{I,J}(W_{q,\varepsilon}^\top)
   &=1+\varepsilon^{22}t_1+o(1).
 \label{eq:kernel-dilution}
\end{align}
These errors tend to zero with $q$ after $\varepsilon$ is fixed.
In particular, the limiting product is
\begin{equation}
 (1+\varepsilon^{13}t_0)(1+\varepsilon^{22}t_1)
 =1+\varepsilon^{13}
      \bigl(t_0+\varepsilon^9t_1+\varepsilon^{22}t_0t_1\bigr)<1
 \label{eq:kernel-product-gap}
\end{equation}
for all sufficiently small positive $\varepsilon$, because the
bracket tends to $t_0<0$. Choose one such $\varepsilon$ and then a
sufficiently large odd prime $q$. The cancellation of $c_0$ in the
coordinate expansion makes the limiting coefficients independent of
the congruence class of $q$. The resulting fixed kernel $W$ is
nonnegative, finite, has mean one, and satisfies
\begin{equation}
 t_{I,J}(W)t_{I,J}(W^\top)<1.
 \label{eq:kernel-oriented-gap}
\end{equation}

\begin{proposition}
\label{prop:negative-kernel}
Assuming Lemma~\ref{lem:rank-layer}, there is a symmetric kernel
$U:\Omega\times\Omega\to[0,1]$ on a finite probability space with
mean $\mu>0$ such that
\begin{equation}
 t(H,U)<\mu^{66}.
 \label{eq:kernel-symmetric-gap}
\end{equation}
\end{proposition}

\begin{proof}
Use the kernel $W:\mathcal X\times\mathcal Y\to[0,\infty)$ just
constructed and set $\Omega=\mathcal X\times\mathcal Y$ with product
probability. Define
\begin{equation}
 V\bigl((x,y),(x',y')\bigr)=W(x,y')W(x',y).
 \label{eq:kernel-crossed-product}
\end{equation}
This kernel is symmetric, nonnegative, and has mean one: the two
factors on the right use independent variables and each has mean
one. In its product over $H$, the factors using $x_i$ for $i\in I$
and $y_j$ for $j\in J$ are independent of those using $y_i$ for
$i\in I$ and $x_j$ for $j\in J$. Therefore
\[
 t(H,V)=t_{I,J}(W)t_{I,J}(W^\top)<1.
\]
Let $L=\max V$, which is finite and at least one, and set $U=V/L$.
Then $0\le U\le1$, $\mu=\E U=1/L>0$, and homogeneity gives
$t(H,U)=L^{-66}t(H,V)<L^{-66}=\mu^{66}$.
\end{proof}

We use the standard independent sampling construction for graphons
\cite[Section~2.6, Lemma~2.4]{lovaszszegedy2006}. The following argument keeps the collision
error explicit and counts all homomorphisms.

\begin{lemma}[Transfer to a finite simple host]
\label{lem:finite-host}
Suppose $U$ is a symmetric $[0,1]$-valued kernel on a finite
probability space, $\mu=\E U>0$, and
\[
 \tau=t(H,U),\qquad \Delta=\mu^{66}-\tau>0.
\]
For any integer
\begin{equation}
 n>\frac{595+66\mu^{66}}{\Delta},
 \label{eq:kernel-host-size}
\end{equation}
there is a finite simple graph $G$ on $n$ vertices with
$t(H,G)<p(G)^{66}$.
\end{lemma}

\begin{proof}
Sample labels $z_1,\ldots,z_n$ independently from the probability
space of $U$. Conditional on these labels, put an edge on each
unordered pair of distinct vertices $\{a,b\}$ independently with
probability $U(z_a,z_b)$. No loops are drawn. This defines a random
finite simple undirected graph $G_n$, and
\begin{equation}
 \E p(G_n)=(1-1/n)\mu.
 \label{eq:kernel-host-mean}
\end{equation}

For every injective map $f:V(H)\to\{1,\ldots,n\}$, its image
vertices have independent labels, and its $66$ source edges use
distinct unordered host pairs. Its probability of being a
homomorphism is exactly $\tau$. Among all $n^{35}$ maps, the
proportion that identify any two source vertices is at most
$\binom{35}{2}/n=595/n$, by the union bound. Each such map has a
homomorphism indicator between zero and one; a collision on a
source edge gives zero, while any repeated host edges are counted
by their actual single Bernoulli trials. No independence claim for
these maps is needed. Comparing their contribution with $\tau$
times their proportion gives
\begin{equation}
 |\E t(H,G_n)-\tau|\le\frac{595}{n}.
 \label{eq:kernel-collision-bound}
\end{equation}
Thus the density still counts all homomorphisms, including every
noninjective one.

Using $(1-1/n)^{66}\ge1-66/n$, Equations~\eqref{eq:kernel-host-size}
and~\eqref{eq:kernel-collision-bound} give
\begin{align*}
 \E t(H,G_n)
 &\le\mu^{66}-\Delta+595/n\\
 &<\mu^{66}-66\mu^{66}/n\\
 &\le\bigl((1-1/n)\mu\bigr)^{66}
   =\bigl(\E p(G_n)\bigr)^{66}
   \le\E\bigl[p(G_n)^{66}\bigr].
\end{align*}
The last inequality is convexity. Hence some realization has
$t(H,G_n)<p(G_n)^{66}$. Such a realization cannot be edgeless.
\end{proof}

Proposition~\ref{prop:negative-kernel} and Lemma~\ref{lem:finite-host}
prove Theorem~\ref{thm:main} conditional only on
Lemma~\ref{lem:rank-layer}. To summarize the order of choices, fix a
sufficiently large even $D$, the $66$ coordinates and $\theta$, then
the scores and a finite $\rho$ satisfying Equation~\eqref{eq:kernel-finite-gap},
then a positive $\varepsilon$ satisfying
Equation~\eqref{eq:kernel-product-gap}, and then a sufficiently large
odd prime $q$. Only after fixing and scaling this finite kernel are
$\mu,\Delta$ and the host size $n$ chosen.
Sections~\ref{sec:transverse}--\ref{sec:tail} establish the required
rank-layer limit.
